\documentclass[10pt,a4paper]{amsart}
\usepackage{amsmath,amssymb,amsthm,graphicx}
\usepackage{verbatim}
\newenvironment{reference}{\comment}{\endcomment}
%\newenvironment{reference}{}{}
\newenvironment{slogan}{\comment}{\endcomment}
\newenvironment{history}{\comment}{\endcomment}

% For commutative diagrams we use Xy-pic
\usepackage[all]{xy}

% We use 2cell for 2-commutative diagrams.
\xyoption{2cell}
% Theorem environments.
%
\theoremstyle{plain}
\newtheorem{theorem}[subsection]{Theorem}
\newtheorem{proposition}[subsection]{Proposition}
\newtheorem{lemma}[subsection]{Lemma}

\theoremstyle{definition}
\newtheorem{definition}[subsection]{Definition}
\newtheorem{example}[subsection]{Example}
\newtheorem{exercise}[subsection]{Exercise}
\newtheorem{situation}[subsection]{Situation}

\theoremstyle{remark}
\newtheorem{remark}[subsection]{Remark}
\newtheorem{remarks}[subsection]{Remarks}

\numberwithin{equation}{subsection}

% Macros
%
\def\lim{\mathop{\mathrm{lim}}\nolimits}
\def\colim{\mathop{\mathrm{colim}}\nolimits}
\def\Spec{\mathop{\mathrm{Spec}}}
\def\Hom{\mathop{\mathrm{Hom}}\nolimits}
\def\Ext{\mathop{\mathrm{Ext}}\nolimits}
\def\SheafHom{\mathop{\mathcal{H}\!\mathit{om}}\nolimits}
\def\SheafExt{\mathop{\mathcal{E}\!\mathit{xt}}\nolimits}
\def\Sch{\mathit{Sch}}
\def\Mor{\mathop{\mathrm{Mor}}\nolimits}
\def\Ob{\mathop{\mathrm{Ob}}\nolimits}
\def\Sh{\mathop{\mathit{Sh}}\nolimits}
\def\NL{\mathop{N\!L}\nolimits}
\def\CH{\mathop{\mathrm{CH}}\nolimits}
\def\proetale{{pro\text{-}\acute{e}tale}}
\def\etale{{\acute{e}tale}}
\def\QCoh{\mathit{QCoh}}
\def\Ker{\mathop{\mathrm{Ker}}}
\def\Im{\mathop{\mathrm{Im}}}
\def\Coker{\mathop{\mathrm{Coker}}}
\def\Coim{\mathop{\mathrm{Coim}}}

% Boxtimes
%
\DeclareMathSymbol{\boxtimes}{\mathbin}{AMSa}{"02}

%
% Macros for moduli stacks/spaces
%
\def\QCohstack{\mathcal{QC}\!\mathit{oh}}
\def\Cohstack{\mathcal{C}\!\mathit{oh}}
\def\Spacesstack{\mathcal{S}\!\mathit{paces}}
\def\Quotfunctor{\mathrm{Quot}}
\def\Hilbfunctor{\mathrm{Hilb}}
\def\Curvesstack{\mathcal{C}\!\mathit{urves}}
\def\Polarizedstack{\mathcal{P}\!\mathit{olarized}}
\def\Complexesstack{\mathcal{C}\!\mathit{omplexes}}
% \Pic is the operator that assigns to X its picard group, usage \Pic(X)
% \Picardstack_{X/B} denotes the Picard stack of X over B
% \Picardfunctor_{X/B} denotes the Picard functor of X over B
\def\Pic{\mathop{\mathrm{Pic}}\nolimits}
\def\Picardstack{\mathcal{P}\!\mathit{ic}}
\def\Picardfunctor{\mathrm{Pic}}
\def\Deformationcategory{\mathcal{D}\!\mathit{ef}}
\usepackage{cleveref}
\crefname{originallemma}{original lemma}{original lemmas}
\Crefname{originallemma}{Original lemma}{Original lemmas}
\makeatletter
\newlabel{lemma-depth-in-ses}{{00LX}{source}{}{sourcepointer-1}{}}
\newlabel{lemma-depth-in-ses@cref}{{[originallemma][0][]00LX}{[1][0][]source}{}{}{}}
\makeatother
\title{The acyclicity lemma: the top cohomology of a bounded complex with depth constraints has depth at least one}
\author{The Stacks Project Authors}
\date{stacks-project:algebra.tex:local-snapshot-20261012; GNU FDL 1.2 or later}
\begin{document}
\maketitle
\noindent\textit{Editorial note.} All delivered mathematical text is the author's own, copied from the pinned Stacks Project chapter file algebra.tex; no statement or proof step was rewritten, reordered or removed.\par\medskip
\noindent\textit{Editorial note.} Preamble provenance, as required for this source: the frozen snapshot stages the chapter's own preamble file beside it (preamble.tex, 217 lines, 8116 bytes, the same shared Stacks-wide preamble the published chow and etale-cohomology cards use); the first line of the pinned chapter file is \textbackslash{}input\{preamble\}, so no preamble text lives inside the chapter itself. The card takes that file's lines 8-18 (the verbatim package and the three \textbackslash{}comment-based environments reference/slogan/history at lines 9-12, section 13.2, together with the Xy-pic setup) and its lines 154 to the runtime end of the file (the native theorem declarations and mathematical macros). Not one character was changed, no package was added, no macro was defined, and the excerpt serves only as preamble, never as a span coordinate.\par\medskip
\noindent\textit{Editorial note.} The same-chapter pointers are located at their \textbackslash{}label lines inside the pinned chapter file; the cross-chapter pointers are located in the project's own tag list staged in the snapshot (sources/stacks-chow/tags.txt), at the verbatim lines carrying their canonical tags.\par\medskip
\noindent\textit{Editorial note.} The counted statement contains a reference environment carrying a citation to the literature. The chapter's own preamble defines that environment as a comment environment (lines 9-12, covered by the preamble spans of this card), so the citation is not typeset and no bibliography is needed; the text is kept verbatim inside the span.\par\medskip
\noindent\textit{Editorial note.} The counted proof refers to the depth lemma of the same chapter by label; it is recorded as a navigation stub pointed at its \textbackslash{}label line, and the argument consuming it stays inside the counted span.\par\medskip
\noindent\textit{Allowed standard prerequisites (editorial).}\begin{itemize}
\item Depth in a short exact sequence. Invoked repeatedly in the counted proof (line 25030 and after). Transfers depth bounds across each of the short exact sequences into which the complex is broken.
\item The depth of the terms of the complex. The hypothesis of the statement, used at the start of the counted proof (lines 25031-25033). Seeds the induction that produces the depth bounds on the intermediate modules.
\end{itemize}\medskip
\ifdefined\hypertarget\hypertarget{sourcepointer-1}{}\fi
\noindent\textit{Source reference.} Original lemma 00LX: sources/stacks-algebra/pinned-original-source.tex, source lines 18122--18122. Navigation pointer for \textbackslash{}ref\{lemma-depth-in-ses\} used inside the counted span (line 25035); the argument that consumes it is retained in the span.\par\medskip
\begin{lemma}[Acyclicity lemma]
\label{lemma-acyclic}
\begin{reference}
\cite[Lemma 1.8]{Peskine-Szpiro}
\end{reference}
Let $R$ be a local Noetherian ring.
Let $0 \to M_e \to M_{e-1} \to \ldots \to M_0$
be a complex of finite $R$-modules.
Assume $\text{depth}(M_i) \geq i$.
Let $i$ be the largest index such that the complex is
not exact at $M_i$. If $i > 0$ then
$\Ker(M_i \to M_{i-1})/\Im(M_{i + 1} \to M_i)$
has depth $\geq 1$.
\end{lemma}

\begin{proof}
Let $H = \Ker(M_i \to M_{i-1})/\Im(M_{i + 1} \to M_i)$ be the
cohomology group in question.
We may break the complex into short exact sequences
$0 \to M_e \to M_{e-1} \to K_{e-2} \to 0$,
$0 \to K_j \to M_j \to K_{j-1} \to 0$, for $i + 2 \leq j \leq e-2 $,
$0 \to K_{i + 1} \to M_{i + 1} \to B_i \to 0$,
$0 \to K_i \to M_i \to M_{i-1}$, and
$0 \to B_i \to K_i \to H \to 0$.
We proceed up through these complexes to
prove the statements about depths, repeatedly using
Lemma \ref{lemma-depth-in-ses}.
First of all, since $\text{depth}(M_e) \geq e$,
and $\text{depth}(M_{e-1}) \geq e-1$ we deduce
that $\text{depth}(K_{e-2}) \geq e - 1$. At this point the
sequences $0 \to K_j \to M_j \to K_{j-1} \to 0$ for $i + 2 \leq j \leq e-2 $
imply similarly that $\text{depth}(K_{j-1}) \geq j$ for
$i + 2 \leq j \leq e-2$. The sequence
$0 \to K_{i + 1} \to M_{i + 1} \to B_i \to 0$
then shows that $\text{depth}(B_i) \geq i + 1$. The sequence
$0 \to K_i \to M_i \to M_{i-1}$ shows that $\text{depth}(K_i) \geq 1$
since $M_i$ has depth $\geq i \geq 1$ by assumption.
The sequence $0 \to B_i \to K_i \to H \to 0$ then
implies the result.
\end{proof}

\begin{proof}
Let $H = \Ker(M_i \to M_{i-1})/\Im(M_{i + 1} \to M_i)$ be the
cohomology group in question.
We may break the complex into short exact sequences
$0 \to M_e \to M_{e-1} \to K_{e-2} \to 0$,
$0 \to K_j \to M_j \to K_{j-1} \to 0$, for $i + 2 \leq j \leq e-2 $,
$0 \to K_{i + 1} \to M_{i + 1} \to B_i \to 0$,
$0 \to K_i \to M_i \to M_{i-1}$, and
$0 \to B_i \to K_i \to H \to 0$.
We proceed up through these complexes to
prove the statements about depths, repeatedly using
Lemma \ref{lemma-depth-in-ses}.
First of all, since $\text{depth}(M_e) \geq e$,
and $\text{depth}(M_{e-1}) \geq e-1$ we deduce
that $\text{depth}(K_{e-2}) \geq e - 1$. At this point the
sequences $0 \to K_j \to M_j \to K_{j-1} \to 0$ for $i + 2 \leq j \leq e-2 $
imply similarly that $\text{depth}(K_{j-1}) \geq j$ for
$i + 2 \leq j \leq e-2$. The sequence
$0 \to K_{i + 1} \to M_{i + 1} \to B_i \to 0$
then shows that $\text{depth}(B_i) \geq i + 1$. The sequence
$0 \to K_i \to M_i \to M_{i-1}$ shows that $\text{depth}(K_i) \geq 1$
since $M_i$ has depth $\geq i \geq 1$ by assumption.
The sequence $0 \to B_i \to K_i \to H \to 0$ then
implies the result.
\end{proof}

\end{document}
